\begin{resumo}[Abstract]
Cardiovascular diseases represent the leading cause of mortality in Brazil and worldwide, motivating the development of computational tools to support clinical diagnosis. This work compares two artificial neural network paradigms, Logistic ADALINE and Multilayer Perceptron (MLP), in the task of heart disease prediction, evaluating each architecture in both custom (NumPy) and library-based (scikit-learn) implementations, under three systematically varied training configurations. Experiments were conducted on two public datasets: the UCI Heart Disease Dataset (918 patients, diagnostic screening) and the Heart Failure Clinical Records Dataset (299 patients, mortality prognosis). A total of 24 independent experiments were assessed not only by overall accuracy, but primarily by the number of False Negatives, the most clinically relevant criterion in medical diagnosis. On Dataset 1, the Sklearn MLP achieved the highest overall accuracy (88.04\%), while the NumPy MLP under configuration C3 (3,000 epochs, $\eta = 0{.}0001$) presented the lowest False Negative count (15), indicating clinical superiority. On Dataset 2, the NumPy MLP under configuration C1 (500 epochs, $\eta = 0{.}01$) offered the best balance between accuracy (78.89\%) and clinical criterion (9 False Negatives). Results show that ADALINE, the simpler linear model, matched MLP performance across several configurations, a phenomenon explained primarily by the limited volume of available training data and, possibly, the predominantly linear nature of the problems under study, a hypothesis consistent with the results but empirically indistinguishable from the data insufficiency hypothesis alone. A complementary experiment with the original continuous data confirmed that variable binarization, adopted with the intent of aligning the representation to official clinical categorizations, is not the determining factor behind ADALINE's competitiveness, and that continuous data consistently improves ADALINE models in both datasets, with gains of up to 13.3 percentage points on Dataset 2, while MLP models do not consistently benefit from continuous representation.

\noindent \textbf{Keywords:} ADALINE. Artificial Neural Networks. False Negative. Heart Disease Diagnosis. Heart Disease Prediction. Machine Learning. Multilayer Perceptron.
\end{resumo}
